\chapter{Conclusion and Future Work} 
\label{ch:conclusion}

This dissertation investigates how volumetric video can evolve from a compelling demonstration into a practical medium for communication over the Internet. 
%
Unlike conventional video, volumetric video represents a dynamic scene in three dimensions and allows viewers to choose their position and viewpoint. 
%
This additional freedom enables immersive conferencing and live streaming, but also introduces substantially higher data rates, stringent latency requirements, and expensive processing. 
%
The systems developed in this dissertation carefully address these challenges at two complementary layers of the delivery pipeline.

\parab{Bandwidth-adaptive Volumetric Video Conferencing.} 
%
\livo demonstrates that full-scene volumetric conferencing can achieve the responsiveness and bandwidth adaptation expected of contemporary 2D communication systems. 
%
It reuses mature 2D video codecs and WebRTC to transport tiled color and depth streams, while introducing depth encoding for large physical spaces, adaptive division of bandwidth between color and depth, and pose-aware view culling. 
%
Evaluations show that \livo sustains approximately 30~fps with around 250~ms of end-to-end latency and improves visual quality over the evaluated alternatives. 
%
These results establish a practical near-term path to immersive two-way communication using commodity RGB-D cameras and existing video infrastructure.

\parab{Frame-rate Compression of Native 3D Representations.} 
%
\gsnfs addresses the codec bottleneck for dynamic Gaussian splats and point clouds. 
%
It reformulates octree construction and RAHT around Morton-ordered, level-synchronous parallelism, keeps the complete codec pipeline on the GPU, and improves attribute compression through color decorrelation. 
%
\gsnfs encodes and decodes evaluated dynamic Gaussian-splat frames within a 30~fps frame budget on desktop GPUs, supports decoding at up to approximately 25~fps on a Jetson Orin, and can also operate as a frame-rate point-cloud codec. 
%
This demonstrates that native 3D compression need not remain the slow component of a volumetric delivery pipeline.

Together, \livo and \gsnfs show two viable and complementary paths toward enabling volumetric video conferencing and live streaming. 
%
\livo builds an end-to-end system from mature 2D codecs, whereas \gsnfs makes native 3D coding fast enough to support the next generation of volumetric representations. 
%
More broadly, the dissertation demonstrates that representation, compression, network adaptation, and computation must be designed jointly if volumetric video is to become a first-class Internet medium.

\section{Future Directions}
\label{sec:future}
%
The systems in this dissertation provide foundations for volumetric communication, but several steps remain before it can become as ubiquitous as 2D video.

\parab{Rate-adaptive Native 3D Delivery.} 
%
\gsnfs can encode at different rate--distortion operating points, but does not yet include a controller that selects those points as network capacity changes. 
%
Coupling its fast encoder to congestion control would enable a native 3D variant of \livo and could support live \gs pipelines driven by emerging feed-forward generators. 
%
Such a system would jointly adapt geometry and appearance while preserving frame-rate encoding and decoding while staying within the network budget.

\parab{Democratizing Volumetric Video.}
%
Future systems should extend \livo beyond two-party sessions and support one-to-many or many-to-many volumetric communication.
%
They have to consider the capabilities of diverse devices such as headsets, smartphones, PCs, and autostereoscopic displays. 
%
This would require device-aware delivery of 3D content such as low-end devices might receive 2D transcoded content while high-end devices might receive native 3D content.
%
Edge processing and personalized views could reduce redundant traffic, while representation and quality could be selected according to each receiver's network, display, and compute capabilities. 
%
Sustained deployment on headsets and mobile GPUs will additionally require careful control of memory traffic, energy, and thermal costs.

\parab{Mobile and Headset Deployment.}
%
The end user of volumetric communication is likely to view content on a headset, mobile GPU, or integrated spatial display rather than a desktop workstation. 
%
\gsnfs's Jetson results indicate that native 3D decoding on such hardware is becoming feasible, while \livo identifies receiver-side reconstruction and rendering costs that must remain within a conferencing budget. 
%
Further work should reduce memory traffic, exploit device-specific accelerators, and coordinate decode quality with display resolution and viewpoint. 
%
Energy consumption and thermal limits will be as important as raw frame rate for sustained sessions.

\parab{Volumetric Video Analytics and Spatial AI.}
%
Volumetric video is not only content for human viewing; it is also a temporally rich record of how people and objects occupy and interact within three-dimensional space. 
%
Future analytics systems can use these streams for 3D detection and tracking, action and interaction understanding, scene-graph construction, and learning spatial world models. 
%
These capabilities can support \textit{spatial AI}: systems that reason about geometry, motion, proximity, and physical properties rather than operating only on pixels or text.
%
Efficient volumetric compression and delivery will make such analytics available across cameras, edge platforms, and cloud services, while analytics can in turn guide which regions and attributes should receive network and compute resources.

\parab{A Three-dimensional Web.}
%
The Web today is largely organized around flat pages and fixed 2D media. 
%
Web in 3D could instead make volumetric scenes and objects first-class, interactive content that users explore from arbitrary viewpoints. 
%
This would enable new Web applications: shoppers could inspect and interact with products in 3D; social platforms could host shared volumetric spaces; sports viewers could choose viewpoints and examine plays spatially; and webpage navigation would be an exploration of content in 3D with audio, visual, and spatial cues rather than text-based 2D monoliths. 
%
Realizing this vision will require progressive 3D delivery, bandwidth adaptation, low-latency interaction, support for heterogeneous devices, and edge-cloud processing. 
%
The conferencing and codec techniques developed in this dissertation provide building blocks for that transition.

These directions expand volumetric video from a communication format into a foundation for interactive media and machine understanding of the physical world.